% -*- mode: context; -*-
\chapter{Trigonometric Ratios Solutions}
\startitemize[n, 1*broad]
\item L.H.S. $= \sqrt{\frac{1- \cos A}{1 + \cos A}}$

   Multiplying and dividing with $1 - \cos A,$ we get

   $= \sqrt{\frac{(1 - \cos A)^2}{1 - \cos^2 A}}$

   $= \sqrt{\left(\frac{1 - \cos A}{\sin A}\right)^2}$

   $= \frac{1 - \cos A}{\sin A} = \csc A - \cot A =$ R.H.S.

\item L.H.S. $\sqrt{1 + \tan^2A + 1 + \cot^2A}$ [$\because \sec^2A = 1 + \tan^2A$ and $\csc^2A = 1 + \cot^2 A$ ]

   $= \sqrt{\tan^2A + 2\tan A\cot A + \cot^2A}$ [$\tan A\cot A = 1$ ]

   $= \sqrt{(\tan A + \cot A)^2} =$ R.H.S.

\item L.H.S $= \left(\frac{1 - \sin^2 A}{\sin A}\right)\left(\frac{1 - \cos^2 A}{\cos A}\right)\left(\frac{\sin^2A +
   \cos^2A}{\sin A \cos A}\right)$

   $= \frac{\cos^2A}{\sin A}\frac{\sin^2 A}{\cos A}\frac{1}{\sin A \cos A} = 1$

\item L.H.S. $= \cos^4A - \sin^4A  + 1= (\cos^2A + \sin^2A)(\cos^2A - \sin^2A) + 1$

   $= \cos^2A - \sin^2A + 1 = 2\cos^2A$

\item L.H.S. $= (\sin A + \cos A)(1 - \sin A\cos A) = (\sin A + \cos A)(\sin^2A + \cos^2A - \sin A\cos A)$

   $= \sin^3A + \cos^3A =$ R.H.S.

\item L.H.S $= \frac{\sin A}{1 + \cos A}+\frac{1 + \cos A}{\sin A} = \frac{\sin^2A + (1 + \cos A)^2}{\sin A(1 + \cos A)}$

   $= \frac{\sin^2A + 1 + 2\cos A + \cos^2A}{\sin A(1 + \cos A)} = \frac{1 + 1 + 2\cos A}{(\sin A)(1 + \cos A))}$

   $= \frac{2}{\sin A} = 2\csc A =$ R.H.S.

\item L.H.S. $= \sin^6A - cos^6A = (\sin^2A + \cos^2A)^3 - 3\cos^4A\sin^2A - 3\cos^2A\sin^4A$

   $= 1 - 3\sin^2A\cos^2A(\cos^2A + \sin^2A) = 1 - 3\sin^2A\cos^2A =$ R.H.S.

\item L.H.S. $= \sqrt{\frac{1 - \sin A}{1 + \sin A}}$

   Multiplying and dividing with $1 - \sin A,$ we get

   $= \sqrt{\frac{(1 - \sin A)^2}{1 - \sin^2A}} = \sqrt{\frac{(1 - \sin A)^2}{\cos^2A}} = \sqrt{\left(\frac{1 - \sin A}{\cos
   A}\right)^2}$

   $= \frac{1 - \sin A}{\cos A} = \sec A - \tan A$

\item L.H.S. $= \frac{\csc A}{\csc A - 1} + \frac{\csc A}{\csc A + 1} = \frac{\frac{1}{\sin A}}{\frac{1}{\sin A} - 1} +
   \frac{\frac{1}{\sin A}}{\frac{1}{\sin A} + 1}$

   $= \frac{1}{1 - \sin A} + \frac{1}{1 + \sin A} = \frac{1}{(1 - \sin A)(1 + \sin A)} = \frac{1}{\cos^2A} = \sec^2A =$
   R.H.S.

\item L.H.S. $= \frac{\csc A}{\tan A + \cot A} = \frac{\frac{1}{\sin A}}{\frac{\sin A}{\cos A} + \frac{\cos A}{\sin A}}$

    $= \frac{\frac{1}{\sin A}}{\frac{\sin^2A + \cos^2A}{\sin A\cos A}} = \frac{\frac{1}{\sin A}}{\frac{1}{\sin A\cos A}} =
    \cos A =$ R.H.S.

\item L.H.S. $= (\sec A + \cos A)(\sec A - \cos A) = \sec^2A - \cos^2A = 1 + \tan^2A - 1 + \sin^2A = \tan^2A + \sin^2A =$
    R.H.S.

\item L.H.S. $= \frac{1}{\tan A + \cot A} = \frac{1}{\frac{\sin A}{\cos A} + \frac{\cos A}{\sin A}} = \frac{1}{\frac{\sin^2A +
    \cos^2A}{\sin A\cos A}} = \sin A\cos A =$ R.H.S.

\item L.H.S $= \frac{1 - \tan A}{1 + \tan A} = \frac{1 - \frac{1}{\cot A}}{1 + \frac{1}{\cot A}} = \frac{\cot A - 1}{\cot A +
    1} =$ R.H.S.

\item L.H.S. $= \frac{1 + \tan^2A}{1 + \cot^2A} = \frac{1 + \frac{\sin^2A}{\cos^2A}}{1 + \frac{\cos^2A}{\sin^A}}$

    $= \frac{\frac{\cos^2A + \sin^2A}{\cos^2A}}{\frac{\sin^2A + \cos^2A}{\sin^2A}} = \frac{\sin^2A}{\cos^2A} =$ R.H.S.

\item L.H.S. $= \frac{\sec A - \tan A}{\sec A + \tan A}$

    Multiplying and dividing with $\sec A - \tan A,$ we get

    $\frac{(\sec A - \tan A)^2}{\sec^2A - \tan^2A} = \sec^2A - 2\tan A\sec A + \tan^A = 1 + \tan^2A - 2\tan A\sec A +
    \tan^2A$

    $= 1 - 2\tan A\sec A + 2\tan^2A =$ R.H.S.

\item L.H.S. $= \frac{1}{\sec A - \tan A}$

    Multiplying and dividing with $\sec A + \tan A,$ we get

    $= \frac{\sec A + \tan A}{\sec^2A = \tan^2A} = \sec A + \tan A =$ R.H.S.

\item L.H.S. $= \frac{\tan A}{1 - \cot A} + \frac{\cot A}{1 - \tan A}$

    $= \frac{\frac{\sin A}{\cos A}}{1 - \frac{\cos A}{\sin A}} + \frac{\frac{\cos A}{\sin A}}{1 - \frac{\sin A}{\cos A}}$

    $= \frac{\sin^2A}{\cos A(\sin A - \cos A)} + \frac{\cos^2A}{\sin A(\cos A - \sin A)}$

    $= \frac{\sin^3A - \cos^3A}{\sin A\cos A(\sin A - \cos A)} = \frac{\sin^2A + \cos^2A + \sin A\cos A}{\sin A\cos A}$

    $= \frac{1 + \sin A \cos A}{\sin A\cos A} = 1 + \csc A\sec A =$ R.H.S.

\item L.H.S. $= \frac{\cos A}{1 - \tan A} + \frac{\sin A}{1 - \cot A} = \frac{\cos A}{1 - \frac{\sin A}{\cos A}} + \frac{\sin
    A}{1 - \frac{\cos A}{\sin A}}$

    $= \frac{\cos^2A}{\cos A - \sin A} + \frac{\sin^2A}{\sin A - \cos A} = \frac{\cos^2A - \sin^2A}{\cos A - \sin A}$

    $= \cos A + \sin A =$ R.H.S.

\item L.H.S. $= (\sin A + \cos A)(\tan A + \cot A) = (\sin A + \cos A)(\frac{\sin A}{\cos A} + \frac{\cos A}{\sin A})$

    $= (\sin A + \cos A)\frac{\sin^2A + \cos^2A}{\sin A\cos A} = \frac{\sin A + \cos A}{\sin A\cos A} = \sec A + \csc A =$
    R.H.S.

\item L.H.S. $= \sec^4A - \sec^2A = (1 + \tan^2A) - 1 - \tan^2A = 1 + 2\tan^2A + \tan^4A - 1 - \tan^2A = \tan^4A + \tan^2A =$
    R.H.S.

\item L.H.S. $= \cot^4A + \cot^2A = (\csc^2A - 1)^1 + \csc^2A - 1 = \csc^4A - 2\csc^2A + 1 + \csc^2A -1 = \csc^4A -
    \csc^2A =$ R.H.S.

\item L.H.S. $= \sqrt{\csc^2A - 1} = \sqrt{\cot^2A} = \cot A = \frac{\cos A}{\sin A} = \cos A \csc A =$ R.H.S.

\item L.H.S. $= \sec^2A\csc^2A = (1 + \tan^2A)(1 + \cot^2A) = 1 + \tan^2A + \cot^2A + \tan^2A\cos^2A = 2 + \tan^2A + \cot^2A
    =$ R.H.S.

\item L.H.S. $= \tan^2A - \sin^2A = \frac{\sin^2A}{\cos^2A} - \sin^2A = \frac{\sin^2A - \sin^2A\cos^2A}{\cos^2A} = \sin^2A(1 -
    \cos^2A)\sec^2A = \sin^4A\sec^2A =$ R.H.S.

\item L.H.S. $= (1 + \cot A - \csc A)(1 + \tan A + \sec A) = \left(1 + \frac{\cos A}{\sin A} - \frac{1}{\sin A}\right)\left(1
    + \frac{\sin A}{\cos A} + \frac{1}{\cos A}\right)$

    $= \frac{(\sin A + \cos A - 1)(\sin A + \cos A + 1)}{\sin A\cos A} = \frac{(\sin A + \cos A)^2 - 1}{\sin A\cos A}$

    $= \frac{\sin^2A + \cos^2A + \sin A\cos A - 1 }{\sin A\cos A} = \frac{1 + 2\sin A\cos A - 1}{\sin A\cos A} = 2=$ R.H.S.

\item L.H.S. $= \frac{\cot A\cos A}{\cot A + \cos A} = \frac{\frac{\cos^2 A}{\sin A}}{\frac{\cos A}{\sin A} + \cos A}$

    $= \frac{\cos^2A}{\sin A}{\frac{\sin A}{\cos A(1 + \sin A)}} = \frac{\cos A}{1 + \sin A}$

    Proceeding in same way for R.H.S. we obtain $\frac{1 - \sin A}{\cos A}$

    $\frac{\cos A}{1 + \sin A} = \frac{\cos A(1 - \sin A)}{1 - \sin^2A} = \frac{1 - \sin A}{\cos A} =$ R.H.S.

\item L.H.S. $= \frac{\cot A + \tan B}{\cot B + \tan A} = \frac{\frac{1}{\tan A} + \tan B}{\frac{1}{\tan B} + \tan A}$

    $\frac{\frac{1 + \tan A\tan B}{\tan A}}{\frac{1 + \tan A\tan B}{\tan B}} = \frac{\tan B}{\tan A} = \cot A\tan B =$ R.H.S.

\item L.H.S. $= \left(\frac{1}{\sec^2 A - \cos^2A} + \frac{1}{\csc^2A - \sin^2A}\right)\cos^2A\sin^2A$

    $= \left(\frac{1}{\frac{1}{\cos^2A} - \cos^2A} + \frac{1}{\frac{1}{\sin^2A} - \sin^2A}\right)\cos^2A\sin^2A$

    $= \left(\frac{\cos^2A}{1 - \cos^4A} + \frac{\sin^2A}{1 - \sin^4}\right)\cos^2A\sin^2A$

    $= \left(\frac{\cos^2A}{(1 - \cos^2A)(1 + \cos^2A)} + \frac{\sin^2A}{(1 - \sin^2A)(1 + \sin^2A)}\right)\cos^2A\sin^2A$

    $= \left(\frac{\cos^2A}{\sin^2A(1 + \cos^2A)} + \frac{\sin^2A}{\cos^2A(1 + \sin^2A)}\right)\cos^2A\sin^2A$

    $= \left(\frac{\cos^4A(1 + \sin^2A) + \sin^4A(1 + \cos^2A)}{\sin^2A\cos^2A(1 + \cos^2A))(1 +
    \sin^2A)}\right)\cos^2A\sin^2A$

    $= \left(\frac{\cos^4A + \sin^4A + \cos^2A\sin^2A(\cos^2A + \sin^2A)}{(1 +\cos^2A)(1 + \sin^2A)}\right)$

    $= \left(\frac{(\cos^2A + \sin^2A)^2 - \cos^2A\sin^2A}{1 + \sin^2A + \cos^2A + \sin^2A\cos^2A}\right)$

    $= \frac{1 - \sin^2A\cos^2A}{2 + \sin^2A\cos^2A} =$ R.H.S.

\item L.H.S. $= \sin^8A - \cos^8A = (\sin^4A - \cos^4A)(\sin^4A + \cos^4A)$

    $= (\sin^2A + \cos^2A)(\sin^2A - \cos^2A)((\sin^2A + \cos^2A)^2 - 2\sin^2A\cos^2A)$

    $= (\sin^2A - \cos^2A)(1 - 2\sin^2A\cos^2A) =$ R.H.S.

\item L.H.S. $= \frac{\cos A\csc A - \sin A\sec A}{\cos A + \sin A} = \frac{\frac{\cos A}{\sin A} - \frac{\sin A}{\cos
    A}}{\cos A + \sin A}$

    $= \frac{\cos^2A - \sin^2A}{\sin A\cos A(\cos A + \sin A)} = \frac{\cos A  \sin A}{\sin A\cos A} = \csc A - \sec A =$
    R.H.S.

\item Given, $\frac{1}{\csc A - \cot A} - \frac{1}{\sin A} = \frac{1}{\sin A} - \frac{1}{\csc A + \cot A}$

    Therefore alternatively we can prove that $\frac{1}{\csc A - \cot A} + \frac{1}{\csc A + \cot A} = \frac{2}{\sin A}$

    L.H.S. $= \frac{1}{\frac{1}{\sin A} - \frac{\cos A}{\sin A}} + \frac{1}{\frac{1}{\sin A} + \frac{\cos A}{\sin A}}$

    $= \frac{\sin A}{1 - \cos A} - \frac{\sin A}{1 + \cos A} = \frac{\sin A(1 + \cos A + 1 - \cos A)}{1 - \cos^2 A} =
    \frac{2\sin A}{\sin^2A} = \frac{2}{\sin A} =$ R.H.S.

\item L.H.S. $= \frac{\tan A + \sec A - 1}{\tan A - \sec A + 1} = \frac{\frac{\sin A + 1 - \cos A}{\cos A}}{\frac{\sin A + \cos
    A - 1)}{\cos A}}$

    $= \frac{1 + \sin A - \cos A}{\sin A + \cos A - 1} = \frac{(1 + \sin A - \cos A)(\sin A + \cos A - 1)}{(\sin A + \cos A -
    1)^2}$

    Solving this yields $= \frac{1 + \sin A}{\cos A} =$ R.H.S.

\item L.H.S. $= (\tan A + \csc B)^2 - (\cot B - \sec A)^2 = \tan^2A + \csc^2B + 2\tan A\csc B - \cot^2B - \sec^2A+ 2\sec
    A\cot B$


    $= (\csc^2B - \cot^2B) - (sec^2A - tan^2A) + 2\tan A\cot B\left(\frac{\csc B}{\cot B} + \frac{\sec A}{\tan A}\right)$

    $= 1 - 1 + 2\tan A\cot B\left(\sec B + \csc A\right)  =$ R.H.S.

\item L.H.S. $= 2\sec^2 A - \sec^4A - 2\csc^2A + \csc^4A = 2(1 + \tan^2A) - (1 + \tan^2A)^2 - 2(1 + \cot^2A) + (1 +
    \cot^2A)^2$

    $= 2 + 2\tan^2A - 1 - 2\tan^2A + \tan^4A - 2 - 2\cot^2A + 1 + 2\cot^2A + \cot^4A = \cot^4A - \tan^4A =$ R.H.S.

\item L.H.S. $= (\sin A + \csc A)^2 + (\cos A + \sec A)^2 = \sin^2A + \csc^2A + 2\sin A\csc A + \cos^2A + \sec^2A + 2\sec
    A\csc A$

    $= (\sin^2A + \cos^2A) + 4 + \csc^2A + \sec^2A = 1 + 4 + 1 + \cot^2A + 1 + \tan^2A = \cot^2A + \tan^2A + 7 =$ R.H.S.

\item L.H.S $= (\csc A + \cot A)(1 - \sin A) - (\sec A + \tan A)(1 - \cos A)$

    $= \left(\frac{1}{\sin A} + \frac{\cos A}{\sin A}\right)(1 - \sin A) - \left(\frac{1}{\cos A} + \frac{\sin A}{\cos
    A}\right)(1 - \cos A)$

    $= \frac{(1 + \cos A)(1 - \sin A)}{\sin A} - \frac{(1 + \sin A)(1 - \cos A)}{\cos A}$

    $= \frac{\cos A(1 + \cos A - \sin A - \cos A\sin A) -\sin A(1 + \sin A - \cos A - \sin A\cos A)}{\sin A\cos A}$

    $= \frac{(1 - \sin A\cos A)(\cos A - \sin A) + (\cos A - \sin A)(\cos A + \sin A)}{\sin A\cos A}$

    $= \frac{(\cos A - \sin A)(1 - \cos A\sin A + \cos A + \sin A)}{\sin A\cos A}$

    $= \left(\frac{(\cos A - \sin A)}{\sin A\cos A}\right)(2 - 1 - \cos A\sin A + \cos A + \sin A)$

    $= (\csc A - \sec A)(2 - (1 - \cos A)(1 - \sin A)) =$ R.H.S.

\item L.H.S $= (1 + \cot A + \tan A)(\sin A - \cos A) = \left(1 + \frac{\csc A}{\sec A} + \frac{\sec A}{\csc
    A}\right)\left(\frac{1}{\csc A} - \frac{1}{\sec A}\right)$

    $= \left(\frac{\sec A\csc A + \sec A | \csc A}{\sec A\csc A}\right)\left(\frac{\sec A - \csc A}{\sec A\csc
    a}\right)$

    $= \frac{\sec^3A - \csc^3A}{\sec^2A\csc^2A} = \frac{\sec A}{\csc^2A} - \frac{\csc A}{\sec^2A} =$ R.H.S.

\item Given, $\frac{1}{\sec A - \tan A} - \frac{1}{\cos A} = \frac{1}{\cos A} - \frac{1}{\sec A + \tan A}$

    So alternatively we can prove that $\frac{1}{\sec A - \tan A} + \frac{1}{\sec A + \tan A} = \frac{2}{\cos A}$

    L.H.S. $= \frac{\sec A + \tan A + \sec A - \tan A}{\sec^2A - tan^2A} = 2\sec A = \frac{2}{\cos A} =$ R.H.S.

\item L.H.S. $= 3(\sin A - \cos A)^4 + 4(\sin^6 A + \cos^6 A) + 6(\sin A + \cos A)^2$

    $= 3[(\sin A - \cos A)^2]^2 + 4[(\sin^2A)^3 + (\cos^2A)^3] + 6(\sin^2A + \cos^2A + 2\cos A\sin A)$

    $= 3[(\sin^2A + \cos^2A - 2\sin A\cos A)]^2 + 4(\sin^2A + \cos^2A)(\sin^4A + \cos^4A - \sin^2A\cos^2A) + 6(1 + 2\cos A\sin
    A)$

    $= 3(1 - 2\sin A\cos A)^2 + 4(\sin^4A + \cos^4A - \sin^2A\cos^2A) + 6(1 + 2\cos A\sin A)$

    $= 3(1 + 4\sin^2A\cos^2A - 4\cos A\sin A) + 4[(\sin^2A + \cos^2A)^2 - 2\sin^2A\cos^2A - \sin^2A\cos^2A] + 6(1 + 2\sin
    A\cos A)$

    $= 3 + 12\sin^2A\cos^2A -12\cos A\sin A + 4(1 - 3\sin^2A\cos^2A) + 6 + 12\sin A\cos A$

    $= 13 =$ R.H.S.

\item L.H.S. $= \sqrt{\frac{1 + \cos A}{1 - \cos A}}$

    Multiplying and dividing with $1 + \cos A,$ we get

    $= \sqrt{\frac{(1 + \cos A)^2}{1 - \cos^2A}} = \sqrt{\left(\frac{1 + \cos A}{\sin A}\right)^2} = \csc A + \cot A =$
    R.H.S.

\item L.H.S. $= \frac{\cos A}{1 + \sin A} + \frac{\cos A}{1 - \sin A}$

    $\frac{\cos A(1 - \sin A) + \cos A(1 + \sin A)}{1 - \sin^2A}$

    $= \frac{2\cos A}{\cos^2A} = 2\sec A =$ R.H.S.

\item L.H.S. $= \frac{\tan A}{\sec A - 1} + \frac{\tan A}{\sec A + 1}$

    $= \frac{\tan A\sec A + \tan A + \tan A\sec A - \tan A}{(\sec^2A - 1)}$

    $= \frac{2\tan A\sec A}{\tan^2 A} = 2\csc A =$ R.H.S.

\item L.H.S. $= \frac{1}{1 - \sin A} - \frac{1}{1 + \sin A}$

    $= \frac{1 + \sin A - 1 + \sin A}{1 - \sin^2A} = \frac{2\sin A}{\cos^2A} = 2\sec A\tan A =$ R.H.S.

\item L.H.S. $= \frac{1 + \tan^2 A}{1 + \cot^2 A} = \frac{1 + \frac{\sin^2A}{\cos^2A}}{1 + \frac{\cos^2A}{\sin^2A}}$

    $= \frac{\frac{\sin^2A + \cos^2A}{\cos^2A}}{\frac{\sin^2A + \cos^2A}{\sin^2A}} = \frac{\sin^2A}{\cos^2A} = \tan^2A$

    R.H.S. $= \left(\frac{1 - \tan A}{1 - \cot A}\right)^2 = \left(\frac{1 - \frac{\sin A}{\cos A}}{1 - \frac{\cos A}{\sin A}}\right)$

    $= \left(\frac{\frac{\cos A - sin A}{\cos A}}{\frac{\sin A - \cos A}{\sin A}}\right)^2$

    $= \frac{\sin^2A}{\cos^2A} = \tan^2A$

    Hence, L.H.S. = R.H.S.

\item L.H.S. $= 1 + \frac{2\tan^2 A}{\cos^2 A} = 1 + \frac{2\sin^2A}{\cos^4A} = \frac{\cos^4A + 2\sin^2A}{\cos^4A}$

    $= \frac{(1 - \sin^2A)^2 + 2\sin^2A}{\cos^4A} = \frac{1 - 2\sin^2A + \sin^4A + 2\sin^2A}{\cos^4A}$

    $= \tan^4A + sec^4A =$ R.H.S.

\item L.H.S. $= (1 - \sin A - \cos A)^2 = 1 - 2\cos A - 2\sin A + 2\sin A\cos A + \sin^2A + \cos^2A$

    $= 2 - 2\cos A - 2\sin A + 2\sin A\cos A = 2(1 - \cos A)(1 - \sin A) =$ R.H.S.

\item L. H. S. $= \frac{\cot A + \csc A - 1}{\cot A - \csc A + 1}$

    $= \frac{\cos A + 1 - \sin A}{\cos A + \sin A - 1}$

    Multiplying and dividing with $\cos A - (1 - \sin A),$ we get

    $= \frac{\cos^2A - 1 + 2\sin A - \sin^2A}{\cos^2A + \sin^2A + 1 - 2\cos A - 2\sin A + 2\sin A\cos A}$

    $= \frac{2\sin A - 2\sin^2A}{2(1 - \cos A)(1 - \sin A)} = \frac{\sin A}{1 - \cos A}$

    Multiplyig numerator and denominator with $1 + \cos A,$ we get

    $= \frac{\sin A(1 + \cos A)}{1 - \cos^2A} = \frac{1 + \cos A}{\sin A} =$ R.H.S.

\item L.H.S. $= (\sin A + \sec A)^2 + (\cos A + \csc A)^2$

    $= \left(\frac{\sin A\cos A + 1}{\cos A}\right)^2 + \left(\frac{\cos A\sin A + 1}{\sin A}\right)^2$

    $= (1 + \sin A\cos A)^2 \left(\frac{1}{\cos^2A} + \frac{1}{\sin^2A}\right)$

    $= (1 + \sin A\cos A)^2\left(\frac{\sin^2A + \cos^2A}{\sin^2A\cos^2A}\right)$

    $= \left(\frac{1 + \sin A\cos A}{\sin A\cos A}\right)^2 = (1 + \sec A\csc A)^2 =$ R.H.S.

\item L.H.S. $= \frac{2\sin A\tan A(1 - \tan A) + 2\sin A\sec^2A}{(1 + \tan A)^2}$

    $= \frac{2\sin A(\tan A - \tan^2A + \sec^2A)}{(1 + \tan A)^2} = \frac{2\sin A(1 + \tan A)}{(1 + \tan A)^2}$

    $= \frac{2\sin A}{1 + \tan A} =$ R.H.S.

\item Given, $2\sin A = 2 - \cos A$

    $\cos A = 2 - 2\sin A,$ squaring both sides we, get $\cos^2A = 4 - 8\sin A + 4\sin^2A$

    $1 - \sin^2A = 4 - 8\sin A + 4\sin^2A \Rightarrow 5\sin^2A - 8\sin A + 3 = 0$

    $5\sin^2A - 5\sin A - 3\sin A + 3 = 0 \Rightarrow \sin A = 1, \frac{3}{5}$

\item Given $8\sin A = 4 + \cos A \Rightarrow 8\sin A - 4 = \cos A$

    Squaring both sides, we get

    $64\sin^2A - 64\sin A + 16 = \cos^2A = 1 - \sin^2A$

    $65\sin^2A - 64\sin A + 15 = 0$

    $65\sin^2A - 39\sin A - 25\sin A + 15 = 0$

    $\sin A = \frac{3}{5}, \frac{5}{13}$

\item Given, $\tan A + \sec A = 1.5$

    $\Rightarrow 1 + \sin A = 1.5 \cos A \Rightarrow 2 + 2\sin A = 3\cos A$

    Squaring both sides, we get

    $4 + 8\sin A + 4\sin^2A = 9 - 9\sin^2A$

    $13\sin^2A + 8\sin A - 5 =0$

    $\sin A = -1, \frac{5}{13}$

\item Given, $\cot A + \csc A = 5 \Rightarrow 1 + \cos A = 5\sin A$

    Squarig both sides, we get

    $1 + 2\cos A + \cos^2A = 25\sin^2A = 25(1 - \cos^2A)$

    $26\cos^2A + 2\cos A - 24 = 0$

    $26\cos^2A + 26\cos A - 24\cos A - 24 = 0$

    $\cos A= -1, \frac{12}{13}$

\item $3\sec^4 A + 8 = 10\sec^2A \Rightarrow 3(sec^2A)^2 + 8 = 10(1 + \tan^2A)$

    $\Rightarrow 3(1 + \tan^2A)^2 + 8 = 10 + 10\tan^2A$

    $3 + 6\tan^2A + 3\tan^4A + 8 = 10 + 10\tan^2A$

    $3\tan^4A - 4\tan^2A + 1 = 0$

    $\tan A = \pm1, \pm\frac{1}{\sqrt{3}}$

\item Given, $\tan^2A + \sec A = 5 \Rightarrow \sec^2A - 1 + \sec A = 5$

    $\sec^2A + \sec A - 6 = 0$

    $\sec A = -3, 2\Rightarrow \cos A = -\frac{1}{3}, \frac{1}{2}$

\item Given $\tan A + \cot A = 2 \Rightarrow \tan A + \frac{1}{\tan A} = 2$

    $\tan^2A - 2\tan A + 1 = 0$

    $\tan A = 1 \Rightarrow \sin A = \frac{1}{\sqrt{2}}$

\item Given, $\sec^2A = 2 + 2\tan A \Rightarrow \tan^2A + 1 = 2 + 2\tan A$

    $\tan^2A - 2\tan A - 1 = 0$

    $\tan A = \frac{2 \pm \sqrt{4 + 4}}{2} = 1 \pm \sqrt{2}$

\item Given, $\tan A = \frac{2x(x + 1)}{2x + 1}$

    $\sin A = \frac{2x(x + 1)}{\sqrt{[2x(x + 1)]^2 + (2x + 1)^2}}$

    $\cos A = \frac{2x + 1}{\sqrt{[2x(x + 1)]^2 + (2x + 1)^2}}$

\item Given, $3\sin A + 5\cos A = 5,$ let $5\sin A - 3\cos A = x$

    Squaring and adding, we get

    $9\sin^2A + 25\cos^2A + 30\sin A\cos A + 25\sin^2A + 9\cos^2A - 30\sin A\cos A = 25 + x^2$

    $9(\sin^2A + \cos^2A) + 25(\cos^2A + \sin^2A) = 25 + x^2$

    $34 = 25 + x^2 \Rightarrow x^2 = 9 \Rightarrow x = \pm 3$

\item Given, $\sec A + \tan A = \sec A - \tan A$

    Multiplying both sides with $\sec A + \tan A,$ we get

    $(\sec A + \tan A)^2 = \sec^2A - \tan^2A = 1$

    $\sec A + \tan A = \pm 1$

    We can prove that $\sec A - \tan A$ to be $\pm 1$ by multiplying given equation with $\sec A - \tan A$

\item Given, $\frac{\cos^4 A}{\cos^2 B} + \frac{\sin^4 A}{\sin^2 B} = 1 = \sin^2A + \cos^2A$

    $\Rightarrow \frac{\cos^4 A}{\cos^2 B} - \cos^2A = \sin^2A - \frac{\sin^4A}{\sin^2B} = 0$

    $\Rightarrow \frac{\cos^2A(\cos^2A - \cos^2B)}{\cos^2B} = \frac{\sin^2A(\sin^2B - \sin^2A)}{\sin^2B}$

    $\Rightarrow \frac{\cos^2A(\cos^2A - \cos^2B)}{\cos^2B} = \frac{\sin^2A(1 - \cos^2B - 1 + \cos^2A)}{\sin^2B}$

    $\Rightarrow (\cos^2A - \cos^2B )\left(\frac{\cos^2A}{\cos^2B} - \frac{\sin^2A}{\sin^2B}\right) = 0$

    When $\cos^2A - \cos^2B = 0, \cos^2A = \cos^2B$

    When $\frac{\cos^2A}{\cos^2B} - \frac{\sin^2A}{\sin^2B} = 0$

    $\cos^2A\sin^2B = \sin^2A\cos^2B \Rightarrow \cos^2A(1 - \cos^2B) = (1 - \cos^2A)\cos^2B$

    $\cos^2A = \cos^2B \Rightarrow \sin^2A = \sin^2B$

    i. L.H.S. $= \sin^4A + \sin^4B = (\sin^2A - \sin^2B)^2 + 2\sin^2A\sin^2B = 2\sin^2A\sin^2B =$ R.H.S.

    ii. L.H.S. $= \frac{\cos^4 B}{\cos^2 A} + \frac{\sin^4 B}{\sin^2 A}$

        $= \frac{\cos^4 B}{\cos^2 B} + \frac{\sin^4 B}{\sin^2 B} = \cos^2B + \sin^2B = 1 =$ R.H.S.

\item Given, $\cos A + \sin A = \sqrt{2}\cos A$

    Squaring both sides

    $1 + 2\sin A\cos A = 2\cos^2A$

    $2 - 2\cos^2A = 2\sin^2A = 1 - 2\sin A\cos A = \sin^2A + \cos^2A - 2\sin A\cos A = (\cos A - \sin A)^2$

    $\cos A - \sin A = \pm \sqrt{2}\sin A$

\item Given, $a\cos A - b\sin A = c,$ and let $a\sin A + b\cos A = x$

    Squaring and adding, we get

    $\Rightarrow a^2\cos^2A + b^2\sin^2A - 2ab\cos A\sin A + a^2\sin^2A + b^2\cos^2A + 2ab\sin A\cos A = c^2 + x^2$

    $\Rightarrow a^2(\cos^2A + \sin^2A) + b^2(\sin^2A + \cos^2A) = c^2 + x^2$

    $\Rightarrow a^2 + b^2 = c^2 + x^2 \Rightarrow x = \pm \sqrt{a^2 + b^2 - c^2}$

\item Given, $1 - \sin A = 1 + \sin A$

    Multiplying both sides by $1 - \sin A,$ we get

    $(1 - \sin A)^2 = 1 - \sin^2A = \cos^2A$

    $1 - \sin A = \pm \cos A$

    Similarly if we multiply with $1 + \sin$ we can prove that

    $1 + \sin A = \pm \cos A$

\item Let us solve these one by one:

    i. Given, $\sin^4A + \sin^2A = 1 \Rightarrow \sin^4A = 1 - \sin^2A = \cos^2$

       $\frac{\sin^4A}{\cos^4A} = sec^2A$

       $\frac{1}{\tan^4A} = cos^2A$

       Also from, $\sin^4A = \cos^2A \Rightarrow \tan^2A = \csc^2A \Rightarrow \frac{1}{\tan^2A} = \sin^2A$

       $\Rightarrow \frac{1}{\tan^4A} + \frac{1}{\tan^2A} = \sin^2A + \cos^2A = 1$

    ii. Given, $\sin^4A + \sin^2A = 1 \Rightarrow \sin^4A = 1 - \sin^2A = \cos^2A$

        $\Rightarrow \frac{\sin^2A}{\cos^2A} = \frac{1}{\sin^2A} \Rightarrow \tan^2A = \csc^2A$

        $\tan^2A = 1 + \cot^2A$

        Multiplying both sides with $\tan^2A,$ we get

        $\tan^4A = \tan^2A + 1$

        $\tan^4A - \tan^2A = 1$

\item Given, $\cos^2 - \sin^2 A = \tan^2 B \Rightarrow \frac{\cos^2A - \sin^2A}{\cos^2A + \sin^2A} = \tan^2B$

    Dividing both numerator and denominator of L.H.S. with $\cos^2A,$ we get

    $\frac{1 - \tan^2A}{1 + \tan^2A} = \tan^2B$

    $1 - \tan^2A - \tan^2B - \tan^2A\tan^2B = 0$

    $1 - \tan^2B = \tan^2A(1 + \tan^2B) \Rightarrow \frac{1 - \tan^2B}{1 + \tan^2B} = \tan^2A$

    $\cos^2B - \sin^2B = \tan^2A \Rightarrow 2\cos^2B - 1 = \tan^2A$

\item We will prove this by induction. Let $\sin A + \csc A = 2,$ thus it is true for $n = 1$

    Squaring both sides $\sin^2A + \csc^2A + 2\sin A\csc A = 2^2 \Rightarrow \sin^2A + \csc^2A = 2$

    Thus, it is true for $n = 2$ as well. Let it be true for $n = m-1$ and $n = m$

    $\sin^mA + \csc^mA = 2$

    Multiplying both sides with $\sin A + \csc A,$ we get

    $\sin^{m + 1}A + \csc^{m + 1}A + \sin^mA\csc A + \csc^mA\sin A = 2^2$

    $\sin^{m + 1}A + \csc^{m + 1}A + \sin^{m - 1}A + \csc^{m - 1}A = 4$

    $\sin^{m + 1}A + \csc^{m + 1} = 4 - 2 = 2$

    Thus, we have proven it by induction.

\item L.H.S. $= \sec A + \tan^3A\csc A = \sec A\left(1 + \tan^3A\frac{\csc A}{\sec A}\right)$

    $=\sec A(1 + \tan^3A\cot a) = \sec A(1 + \tan^2A) = \sec^3A = (1 + 1 - e^2)\frac{3}{2} = (2 - e^2)\frac{3}{2}$

\item Cross multiplying given equations, we have

    $\frac{\sec A}{br - qc} = \frac{\tan A}{pc - qr} = \frac{1}{qa - pb}$

    $\therefore \sec^2A - \tan^2A = 1, \left(\frac{br - qc}{aq - pb}\right)^2 - \left(\frac{pc - qa}{aq - pb}\right)^2 = 1$

    $(br - qa)^2 - (pc - ar)^2 = (qa - pb)^2$

\item Given, $\csc A - \sin A = m \Rightarrow \frac{1}{\sin A } - \sin A = m$

    $\Rightarrow \frac{1 - \sin^2A}{\sin A} = m \Rightarrow \frac{\cos^2A}{\sin A} = m$

    Also given, $\sec A - \cos A = n \Rightarrow \frac{1}{\cos A} - \cos A = n$

    $\Rightarrow \frac{1 - \cos^2A}{\cos A} = n \Rightarrow \frac{\sin^2A}{\cos A} = n$

    We have $\sin A = \frac{\cos^2A}{m},$ putting this in derived equation

    $\cos^3A = m^2n$

    $\therefore \sin A = \frac{(m^2n)^{\frac{2}{3}}}{m} = (mn^2)\frac{1}{3}$

    Thus, $\sin^2A + \cos^2A = 1$ gives us $(m^2n)^{\frac{2}{3}} + (mn^2)^{\frac{2}{3}} = 1$

\item Given, $\sec^2A = \frac{4xy}{(x + y)^2}$

    $\because \sec^2A \geq 1 \therefore \frac{4xy}{(x + y)^2}\geq 1$

    $\Rightarrow (x + y)^2 \leq 4xy \Rightarrow (x - y)^2 \leq 0$

    But for real $x$ and $y, (x - y)^2\nless 0$

    $\therefore (x - y)^2 = 0 \therefore x - y$

    Also, $x + y \neq 0 \Rightarrow x\neq 0, y \neq 0$

\item Given, $\sin A = x + \frac{1}{x},$ squaring we get

    $\sin^2A = x^2 + \frac{1}{x^2} + 2 \geq 2$

    which is not possible since $\sin A \leq 1$

\item Given, $\sec A - \tan A = p$

    $\Rightarrow 1 - \sin A = p\cos A$

    Squaring, we obtain

    $1 + \sin^2A - 2\sin A = p^2\cos^2A = p^2(1 - \sin^2A)$

    $(1 + p^)\sin^2A - 2\sin A + 1 - p^2 = 0$

    $\sin A = \frac{1 \pm p^2}{1 + p^2}$

    Now $tan$ and $\sec A$ can be easily found.

\item $\sec A + \tan A = \frac{4p^2 + 1}{4p} \pm \sqrt{\sec^2A - 1}$

    $= 4p \pm \sqrt{\left(\frac{(4p^2 + 1)^2}{16p^2} - 1\right)}$

    $= 2p$ or $\frac{1}{2p}$

\item Given, $\frac{\sin A}{\sin B} = p, \frac{\cos A}{\cos B} = q$

    Squaring, we get

    $\frac{\sin^2A}{\sin^2B} = p^2, \frac{\cos^2A}{\cos^2B}= q^2$

    $\frac{\sin^2A - \sin^2B}{\sin^2B} = p^2 - 1, \frac{\cos^2B - \cos^2A}{\cos^2B} = 1 - q^2$

    $\frac{\sin^2A - \sin^2B}{\sin^2B} = p^2 - 1, \frac{\sin^2A - \sin^2B}{\cos^2B} = 1 - q^2$

    Dividing, we obtain

    $\tan^2B = \pm \sqrt{\frac{1 - q^2}{p^2 - 1}}$

    Dividing original equations

    $\frac{\tan A}{\tan B} = \pm \frac{p}{q}\sqrt{\frac{1 - q^2}{p^2 - 1}}$

\item $\sin A = \sqrt{2}\sin B, \frac{\sin A.\cos B}{\cos A. \sin B} = \sqrt{3}$

    Substituting for $\sin A$

    $\frac{\sqrt{2}\sin B\cos B}{\sqrt{1 - 2\sin^2A}\sin B}  = \sqrt{3}$

    Squaring, we get

    $2\cos^2 A = 3(1 - 2\sin^2A) \Rightarrow 4\sin^2A - 1 = 0$

    $\Rightarrow \sin A = \pm \frac{1}{2}, \Rightarrow A = \pm 45^\circ$

    Thus, $B = \pm 30^\circ$

\item Given, $\tan A + \cot A = 2, 1 + \tan^2A = 2\tan A \Rightarrow (1 - \tan A)^2 = 0$

    $\tan A = 1\Rightarrow \sin A = \pm\frac{1}{\sqrt{2}}$

\item $m^2 - n^2 = (\tan A +\sin A)^2 - (\tan A - \sin A)^2 = 4\tan A\sin A$

    $4\sqrt{mn} = 4\sqrt{\tan^2A - \sin^2A} = 4\sqrt{\frac{\sin^2A}{\cos^2A} - \sin^2A}$

    $= 4\sqrt{\frac{\sin^2A(1 - \cos^2A)}{\cos^2A}} = 4\sqrt{\frac{\sin^4A}{\cos^2A}}$

    $= 4\tan A\sin A$

\item $m^2 - 1 = 2\sin A\cos A \Rightarrow n(m^2 - 1) = 2\sin A\cos A(\sec A + \csc A)$

    $= 2\sin A + 2\cos A = 2m$

\item Given, $x\sin^3 A + y\cos^3 A = \sin A\cos A$

    $(x\sin A)\sin^2A + (y\cos A)\cos^2A = \sin A\cos A$

    Frpm other equation $x\sin A = y\cos A,$ substituting this in above equation

    $(x\sin A)\sin^2A + (x\sin A)\cos^2a = \sin A\cos A$

    $x\sin A = \sin A\cos A$

    $x = \cos A \therefore y = \sin A$

    Thus, $x^2 + y^2 = 1$

\item Given, $\sin^2A = \frac{(x + y)^2}{4xy}$

    $\sin^2A \leq 1 \Rightarrow (x + y)^2 \leq 4xy \Rightarrow (x - y)^2 \leq 0$

    But for real $x$ and $y$ $(x - y)^2\nless 0$

    $\Rightarrow x = y$

    Also, $xy \neq 0, x, y\neq 0$
\stopitemize

\subject{Challenging Problems}
\startitemize[n, 1*broad]
\item Given, $\sin A + \sin^3 A = 1 - \sin^2 A = \cos^2 A,$ i.e. $\sin A(1 + \sin^2 A) = \cos^2 A$.

    Squaring, $\sin^2 A(1 + \sin^2 A)^2 = \cos^4 A \Rightarrow (1 - \cos^2 A)(2 - \cos^2 A)^2 = \cos^4 A$

    $\Rightarrow 4 - 8\cos^2 A + 5\cos^4 A - \cos^6 A = \cos^4 A$

    $\Rightarrow \cos^6 A - 4\cos^4 A + 8\cos^2 A = 4$

\item Given, $\cos A = 1 - \cos^2 A = \sin^2 A$. Hence $\sin^4 A + \sin^2 A = \cos^2 A + \cos A = 1$.

    $\sin^{12} A + 3\sin^{10} A + 3\sin^8 A + \sin^6 A = (\sin^4 A + \sin^2 A)^3 = 1$

    $2\sin^4 A + 2\sin^2 A - 2 = 2(\sin^4 A + \sin^2 A) - 2 = 0$

    Adding, L.H.S. $= 1$.

\item $\sin B = \frac{\sin A}{m}$ and $\tan B = \frac{\tan A}{n} \Rightarrow \cos B = \frac{\sin B}{\tan B} = \frac{n\cos A}{m}$

    $\sin^2 B + \cos^2 B = 1 \Rightarrow \sin^2 A + n^2\cos^2 A = m^2$

    $\Rightarrow 1 - \cos^2 A + n^2\cos^2 A = m^2 \Rightarrow \cos^2 A = \frac{m^2 - 1}{n^2 - 1}$

\item $(x\sin A - y\cos A)^2 + (x\cos A + y\sin A)^2 = x^2 + y^2$

    The first equation gives $(x\cos A + y\sin A)^2 = 0,$ so $\tan A = -\frac{x}{y}$

    $\Rightarrow \sin^2 A = \frac{x^2}{x^2 + y^2}, \cos^2 A = \frac{y^2}{x^2 + y^2}$

    Substituting in the second equation, $\frac{y^2}{a^2} + \frac{x^2}{b^2} = 1$

\item $\sin A\cos A = \frac{a^2 - 1}{2}$

    $b = (\sin A + \cos A)(1 - \sin A\cos A) = a\left(1 - \frac{a^2 - 1}{2}\right) = \frac{a(3 - a^2)}{2}$

    Required result: $a^3 - 3a + 2b = 0$

\item $a = \frac{1}{\sin A\cos A}, b = \frac{1 - \cos^2 A}{\cos A} = \frac{\sin^2 A}{\cos A}$

    $\frac{b}{a} = \sin^3 A \Rightarrow \sin A = \left(\frac{b}{a}\right)^{\frac{1}{3}}, \cos A = \frac{\sin^2 A}{b} =
    \frac{1}{a^{\frac{2}{3}}b^{\frac{1}{3}}}$

    $\sin^2 A + \cos^2 A = 1 \Rightarrow \frac{b^{\frac{2}{3}}}{a^{\frac{2}{3}}} + \frac{1}{a^{\frac{4}{3}}b^{\frac{2}{3}}} = 1$

    Multiplying by $a^{\frac{4}{3}}b^{\frac{2}{3}},$ we get $(a^2b)^{\frac{2}{3}} - (ab^2)^{\frac{2}{3}} = 1$

\item Let $t = \sin A\cos A = \frac{b^2 - 1}{2}$

    $a = \frac{\sin^2 A - \cos^2 A}{t} \Rightarrow a^2 = \frac{(\sin^2 A + \cos^2 A)^2 - 4t^2}{t^2} = \frac{1}{t^2} - 4$

    $\Rightarrow (a^2 + 4)t^2 = 1 \Rightarrow (a^2 + 4)(b^2 - 1)^2 = 4$

\item $\sin A\cos A = \frac{a^2 - 1}{2}$

    $b = (\sin^2 A + \cos^2 A)^3 - 3\sin^2 A\cos^2 A(\sin^2 A + \cos^2 A) = 1 - \frac{3(a^2 - 1)^2}{4}$

    Required result: $3(a^2 - 1)^2 = 4(1 - b)$

\item Put $\cos^2 A = 1 - \sin^2 A$ and clear fractions:

    $(a + b)\left[b\sin^4 A + a(1 - \sin^2 A)^2\right] = ab$

    $\Rightarrow (a + b)^2\sin^4 A - 2a(a + b)\sin^2 A + a^2 = 0 \Rightarrow \left[(a + b)\sin^2 A - a\right]^2 = 0$

    $\therefore \sin^2 A = \frac{a}{a + b}, \cos^2 A = \frac{b}{a + b}$

    $\frac{\sin^8 A}{a^3} + \frac{\cos^8 A}{b^3} = \frac{a}{(a + b)^4} + \frac{b}{(a + b)^4} = \frac{1}{(a + b)^3}$

\item From the first, $m = \frac{a\tan^2 A + b}{1 + \tan^2 A} \Rightarrow \tan^2 A = \frac{m - b}{a - m}$

    From the second, $n = \frac{b\tan^2 B + a}{1 + \tan^2 B} \Rightarrow \tan^2 B = \frac{n - a}{b - n}$

    $a^2\tan^2 A = b^2\tan^2 B \Rightarrow a^2(m - b)(b - n) = b^2(n - a)(a - m)$

    Expanding, and dividing by $(a - b)$ (assuming $a \neq b$), we get $ab(m + n) = mn(a + b)$

    i.e. $\frac{1}{a} + \frac{1}{b} = \frac{1}{m} + \frac{1}{n}$

\item Let each expression be $k$. Then $x = \frac{k}{\sin A}, y = \frac{k}{\cos A}$

    $z = \frac{k(1 - \tan^2 A)}{2\tan A} = \frac{k(\cos^2 A - \sin^2 A)}{2\sin A\cos A}$

    $x^2 + y^2 = \frac{k^2}{\sin^2 A\cos^2 A}, x^2 - y^2 = \frac{k^2(\cos^2 A - \sin^2 A)}{\sin^2 A\cos^2 A}$

    $4z^2(x^2 + y^2) = \frac{k^4(\cos^2 A - \sin^2 A)^2}{\sin^4 A\cos^4 A} = (x^2 - y^2)^2$

\item $(\sec A - \sec B\sec C)^2 = \tan^2 B\tan^2 C = (\sec^2 B - 1)(\sec^2 C - 1)$

    $\Rightarrow \sec^2 A + \sec^2 B + \sec^2 C - 2\sec A\sec B\sec C = 1$

    This relation is symmetric in $A, B, C$. Rearranging it,

    $(\sec B - \sec A\sec C)^2 = (\sec^2 A - 1)(\sec^2 C - 1) = \tan^2 A\tan^2 C$

    $\Rightarrow \sec B = \sec A\sec C \pm \tan A\tan C$

\item Let $x = \tan^2 A, y = \tan^2 B, z = \tan^2 C$. Then $\sin^2 A = \frac{x}{1 + x}$ etc.

    $\frac{x}{1 + x} + \frac{y}{1 + y} + \frac{z}{1 + z} = 1 \Leftrightarrow \sum x(1 + y)(1 + z) = (1 + x)(1 + y)(1 + z)$

    $\Leftrightarrow (x + y + z) + 2(xy + yz + zx) + 3xyz = 1 + (x + y + z) + (xy + yz + zx) + xyz$

    $\Leftrightarrow xy + yz + zx + 2xyz = 1,$ which is given.

\item From the first, $\sin A(1 - a\cos B) = a\sin B\cos A \Rightarrow \sin A = a(\sin A\cos B + \cos A\sin B)$

    From the second, $\sin B(1 - b\cos A) = b\sin A\cos B \Rightarrow \sin B = b(\sin A\cos B + \cos A\sin B)$

    Dividing, $\frac{\sin A}{\sin B} = \frac{a}{b}$

\item $\frac{1 - \cos A}{1 + \cos A} = \frac{b + c - a}{b + c + a}$

    Similarly the others are $\frac{c + a - b}{a + b + c}$ and $\frac{a + b - c}{a + b + c}$

    Sum $= \frac{(b + c - a) + (c + a - b) + (a + b - c)}{a + b + c} = 1$

\item (i) $1 - \cos^2 A = \frac{(1 - e\cos B)^2 - (\cos B - e)^2}{(1 - e\cos B)^2} = \frac{(1 - e^2)(1 - \cos^2 B)}{(1 - e\cos
    B)^2}$

    Hence $\sin A = \frac{\sqrt{1 - e^2}\sin B}{1 - e\cos B}$ (both sides positive).

    (ii) $1 - \cos A = \frac{(1 + e)(1 - \cos B)}{1 - e\cos B}, 1 + \cos A = \frac{(1 - e)(1 + \cos B)}{1 - e\cos B}$

    Dividing gives the result.

\item Let $u = \tan A + \cot A \geq 2$. Expression $= u^2 - 2 - 6u = (u - 3)^2 - 11$

    Least value is $-11,$ when $u = 3$ i.e. $\tan^2 A - 3\tan A + 1 = 0,$ $\tan A = \frac{3 \pm \sqrt{5}}{2}$.

\item $\cos^2 A = \sin A\tan A = \frac{\sin^2 A}{\cos A} \Rightarrow \cos^3 A = \sin^2 A$

    $\cot^2 A = \frac{\cos^2 A}{\cos^3 A} = \frac{1}{\cos A}, \cot^6 A = \frac{1}{\cos^3 A}$

    $\cot^6 A - \cot^2 A = \frac{1 - \cos^2 A}{\cos^3 A} = \frac{\sin^2 A}{\sin^2 A} = 1$

\item $\sec^2 A = 1 + \tan^2 A = \frac{(\sin B + \cos B)^2 + (\sin B - \cos B)^2}{(\sin B + \cos B)^2} = \frac{2}{(\sin B + \cos
    B)^2}$

    $\Rightarrow (\sin B + \cos B)^2 = 2\cos^2 A \Rightarrow \sin B + \cos B = \pm\sqrt{2}\cos A$

\item Given, $\cos A\cos B - \sin A\sin B = 1$. Now

    $(\cos A\cos B - \sin A\sin B)^2 + (\sin A\cos B + \cos A\sin B)^2 = (\sin^2 A + \cos^2 A)(\sin^2 B + \cos^2 B) = 1$

    $\Rightarrow \sin A\cos B + \cos A\sin B = 0$

    Dividing by $\sin A\cos B,$ $1 + \cot A\tan B = 0$

\item Let $p = m^2 + 2mn, h = m^2 + 2mn + 2n^2$.

    $h^2 - p^2 = (h - p)(h + p) = 2n^2(2m^2 + 4mn + 2n^2) = 4n^2(m + n)^2$

    $\cos A = \frac{2n(m + n)}{h} \Rightarrow \tan A = \frac{p}{2n(m + n)} = \frac{m^2 + 2mn}{2mn + 2n^2}$

\item L.H.S. $= \frac{(a + b) + (a - b)}{\sqrt{a^2 - b^2}} = \frac{2a}{\sqrt{a^2 - b^2}} = \frac{2}{\sqrt{1 - \tan^2 A}}$

    $= \frac{2}{\sqrt{\frac{\cos^2 A - \sin^2 A}{\cos^2 A}}} = \frac{2\cos A}{\sqrt{\cos^2 A - \sin^2 A}}$

\item Squaring, $9\tan^2 A + 24\tan A + 16 = 25(1 + \tan^2 A)$

    $\Rightarrow 16\tan^2 A - 24\tan A + 9 = 0 \Rightarrow (4\tan A - 3)^2 = 0 \Rightarrow \tan A = \frac{3}{4}$

    Then $5\sec A = \frac{25}{4},$ $\sec A = \frac{5}{4}$ (checks). $\sin A = \frac{3}{5}, \cos A = \frac{4}{5}$

\item $\sec^2 A + {\rm cosec}^2 A = \frac{1}{\sin^2 A\cos^2 A} = \frac{16}{3} \Rightarrow \sin A\cos A = \frac{\sqrt{3}}{4}$

    $(\sin A + \cos A)^2 = 1 + \frac{\sqrt{3}}{2} \Rightarrow \sin A + \cos A = \frac{\sqrt{3} + 1}{2}$

    $\sin A, \cos A$ are roots of $x^2 - \frac{\sqrt{3} + 1}{2}x + \frac{\sqrt{3}}{4} = 0,$ i.e. $\frac{1}{2},
    \frac{\sqrt{3}}{2}$

    $A = 30^\circ$ or $60^\circ$

\item Let $p = \sin A + \cos A,$ $\sin A\cos A = \frac{p^2 - 1}{2}$. Given, $p = 2\sqrt{2}\sin A\cos A = \sqrt{2}(p^2 - 1)$

    $\sqrt{2}p^2 - p - \sqrt{2} = 0 \Rightarrow p = \sqrt{2}$ or $-\frac{1}{\sqrt{2}}$. For acute $A, p > 0,$ so $p = \sqrt{2}$

    $\sin A\cos A = \frac{1}{2} \Rightarrow (\sin A - \cos A)^2 = 0 \Rightarrow A = 45^\circ$

\item $\sec^2 A = 1 + (2 - \sqrt{3})^2 = 8 - 4\sqrt{3} = (\sqrt{6} - \sqrt{2})^2$

    $\cos A = \frac{1}{\sqrt{6} - \sqrt{2}} = \frac{\sqrt{6} + \sqrt{2}}{4}$

    $\sin A = \tan A\cos A = \frac{(2 - \sqrt{3})(\sqrt{6} + \sqrt{2})}{4} = \frac{\sqrt{6} - \sqrt{2}}{4}$

\item $\sec^2 A = 1 + (\sqrt{2} - 1)^2 = 4 - 2\sqrt{2} \Rightarrow \cos^2 A = \frac{1}{4 - 2\sqrt{2}} = \frac{2 + \sqrt{2}}{4}$

    $\cos A = \frac{\sqrt{2 + \sqrt{2}}}{2}, \sin^2 A = 1 - \cos^2 A = \frac{2 - \sqrt{2}}{4}, \sin A = \frac{\sqrt{2 -
    \sqrt{2}}}{2}$

\item $\sin^2 A = \frac{6 - 2\sqrt{5}}{16} = \frac{3 - \sqrt{5}}{8}, \cos^2 A = \frac{5 + \sqrt{5}}{8} \Rightarrow \cos A =
    \frac{\sqrt{10 + 2\sqrt{5}}}{4}$

    $\tan^2 A = \frac{3 - \sqrt{5}}{5 + \sqrt{5}} = \frac{(3 - \sqrt{5})(5 - \sqrt{5})}{20} = \frac{20 - 8\sqrt{5}}{20} = 1 -
    \frac{2}{\sqrt{5}}$

\item Let $p = \sin A + \cos A$. Then $(1 + \sin A)(1 + \cos A) = 1 + p + \frac{p^2 - 1}{2} = \frac{(p + 1)^2}{2}$

    $\frac{(p + 1)^2}{2} = \frac{5}{4} \Rightarrow p = -1 \pm \sqrt{\frac{5}{2}}$. Since $p^2 \leq 2,$ $p = \sqrt{\frac{5}{2}} - 1$

    $(1 - \sin A)(1 - \cos A) = \frac{(p - 1)^2}{2} = \frac{\left(\sqrt{\frac{5}{2}} - 2\right)^2}{2} = \frac{13}{4} - \sqrt{10}$

\item $\tan^2 A + \cot^2 A = 16 - 2 = 14$

    $\tan^3 A + \cot^3 A = 4^3 - 3\cdot 4 = 52$

    $\sin A\cos A = \frac{1}{\tan A + \cot A} = \frac{1}{4},$ so $\sin^4 A + \cos^4 A = 1 - 2\cdot\frac{1}{16} = \frac{7}{8}$

    $\sin^6 A + \cos^6 A = 1 - 3\cdot\frac{1}{16} = \frac{13}{16}$

\item Let $t = \sin A\cos A = \frac{p^2 - 1}{2}$. $\sin^4 A + \cos^4 A = 1 - 2t^2$

    $\sin^8 A + \cos^8 A = (1 - 2t^2)^2 - 2t^4 = 1 - 4t^2 + 2t^4$

    $= 1 - (p^2 - 1)^2 + \frac{(p^2 - 1)^4}{8}$

\item Let $x = \sin^2 A, y = \cos^2 A,$ $x + y = 1,$ $xy = t$.

    $x^5 + y^5 = (x + y)^5 - 5xy(x + y)^3 + 5x^2y^2(x + y) = 1 - 5t + 5t^2$

    [since $x^5 + y^5 = (x^2 + y^2)(x^3 + y^3) - x^2y^2(x + y),$ $x^2 + y^2 = 1 - 2t,$ $x^3 + y^3 = 1 - 3t$]

    $= 1 - 5\sin^2 A\cos^2 A + 5\sin^4 A\cos^4 A$

\item $\sin^4 A + \cos^4 A - 1 = -2\sin^2 A\cos^2 A$

    $\sin^6 A + \cos^6 A - 1 = -3\sin^2 A\cos^2 A$

    Ratio $= \frac{2}{3},$ independent of $A$.

\item L.H.S. $= 3(\sin^4 A + \cos^4 A) - 2(\sin^6 A + \cos^6 A)$

    $= 3(1 - 2\sin^2 A\cos^2 A) - 2(1 - 3\sin^2 A\cos^2 A) = 1$

\item $\frac{\tan^3 A}{\sec^2 A} = \frac{\sin^3 A}{\cos A}, \frac{\cot^3 A}{{\rm cosec}^2 A} = \frac{\cos^3 A}{\sin A}$

    Sum $= \frac{\sin^4 A + \cos^4 A}{\sin A\cos A} = \frac{1 - 2\sin^2 A\cos^2 A}{\sin A\cos A}$

    $= \sec A\,{\rm cosec} A - 2\sin A\cos A$

\item $\frac{\tan^2 A}{\tan^2 A - 1} = \frac{\sin^2 A}{\sin^2 A - \cos^2 A}$

    $\frac{{\rm cosec}^2 A}{\sec^2 A - {\rm cosec}^2 A} = \frac{\frac{1}{\sin^2 A}}{\frac{\sin^2 A - \cos^2 A}{\sin^2 A\cos^2 A}} =
    \frac{\cos^2 A}{\sin^2 A - \cos^2 A}$

    Sum $= \frac{\sin^2 A + \cos^2 A}{\sin^2 A - \cos^2 A} = \frac{1}{\sin^2 A - \cos^2 A}$

\item $\frac{\cos^2 A}{1 - \tan A} = \frac{\cos^3 A}{\cos A - \sin A}$

    L.H.S. $= \frac{\cos^3 A - \sin^3 A}{\cos A - \sin A} = \cos^2 A + \sin A\cos A + \sin^2 A = 1 + \sin A\cos A$

\item L.H.S. $= \frac{\sin A - \cos A}{\sin A\cos A}\cdot\frac{\sin A\cos A + \sin^2 A + \cos^2 A}{\sin A\cos A}$

    $= \frac{(\sin A - \cos A)(1 + \sin A\cos A)}{\sin^2 A\cos^2 A} = \frac{\sin^3 A - \cos^3 A}{\sin^2 A\cos^2 A}$

    $= \frac{\sin A}{\cos^2 A} - \frac{\cos A}{\sin^2 A} = \tan A\sec A - \cot A\,{\rm cosec} A$

\item L.H.S. $= \frac{\cos^2 A}{\sin^2 A}\cdot\frac{1 - \cos A}{\cos A(1 + \sin A)} = \frac{\cos A(1 - \cos A)}{(1 - \cos^2 A)(1
    + \sin A)} = \frac{\cos A}{(1 + \cos A)(1 + \sin A)}$

    R.H.S. $= \frac{1}{\cos^2 A}\cdot\frac{(1 - \sin A)\cos A}{1 + \cos A} = \frac{1 - \sin A}{\cos A(1 + \cos A)}$

    $= \frac{1 - \sin^2 A}{\cos A(1 + \cos A)(1 + \sin A)} = \frac{\cos A}{(1 + \cos A)(1 + \sin A)}$

\item $\frac{\sin A}{\sec A + \tan A - 1} = \frac{\sin A\cos A}{1 + \sin A - \cos A},$ $\frac{\cos A}{{\rm cosec} A + \cot A - 1} =
    \frac{\sin A\cos A}{1 + \cos A - \sin A}$

    Sum $= \sin A\cos A\cdot\frac{2}{1 - (\sin A - \cos A)^2} = \frac{2\sin A\cos A}{2\sin A\cos A} = 1$

\item $\frac{1 + \sin A - \cos A}{1 + \sin A + \cos A} = \frac{1 - \cos A}{\sin A},$ since cross-multiplication gives

    $\sin A + \sin^2 A - \sin A\cos A = 1 + \sin A - \sin A\cos A - \cos^2 A,$ i.e. $\sin^2 A = 1 - \cos^2 A$

    $\therefore$ L.H.S. $= \frac{(1 - \cos A)^2}{\sin^2 A} = \frac{(1 - \cos A)^2}{(1 - \cos A)(1 + \cos A)} = \frac{1 - \cos A}{1 +
    \cos A}$

\item L.H.S. $= 1 + \tan^2 A\tan^2 B + \tan^2 A + \tan^2 B$ (the cross terms $\pm 2\tan A\tan B$ cancel)

    $= (1 + \tan^2 A)(1 + \tan^2 B) = \sec^2 A\sec^2 B$

\item Let $t = \tan A > 0$. $f = \frac{t^2 - t + 1}{t^2 + t + 1} = 1 - \frac{2t}{t^2 + t + 1} = 1 - \frac{2}{t + \frac{1}{t} + 1}$

    $t + \frac{1}{t}$ takes every value in $[2, \infty),$ so $\frac{2}{t + \frac{1}{t} + 1}$ takes every value in $\left(0,
    \frac{2}{3}\right]$

    Set of values: $\frac{1}{3} \leq f < 1,$ least value $\frac{1}{3}$ at $A = 45^\circ$.

\item Expanding, cross terms $\pm 2ab\sin A\cos A$ cancel, leaving $(a^2 + b^2)(\sin^2 A + \cos^2 A) = a^2 + b^2$

    With $a = 5, b = 12$: $(5\sin A + 12\cos A)^2 = 169 - (5\cos A - 12\sin A)^2 \leq 169$

    Greatest value $13,$ attained when $\tan A = \frac{5}{12}$.

\item $a^2\sec^2 A + b^2{\rm cosec}^2 A = a^2 + b^2 + a^2\tan^2 A + b^2\cot^2 A$

    $\geq a^2 + b^2 + 2ab = (a + b)^2$ (A.M. $\geq$ G.M.)

    Least value $(a + b)^2,$ when $\tan^2 A = \frac{b}{a}$.

\item Let $\sin A = s, \cos A = c$. By A.M. $\geq$ G.M. with weights,

    $\frac{a}{c} + \frac{a}{c} + k c^2 \geq 3(a^2k)^{\frac{1}{3}}$ and $\frac{b}{s} + \frac{b}{s} + k s^2 \geq 3(b^2k)^{\frac{1}{3}}$
    for any $k > 0$

    Adding: $2(a\sec A + b\,{\rm cosec} A) + k \geq 3k^{\frac{1}{3}}(a^{\frac{2}{3}} + b^{\frac{2}{3}})$

    Choose $k = (a^{\frac{2}{3}} + b^{\frac{2}{3}})^{\frac{3}{2}}$: $a\sec A + b\,{\rm cosec} A \geq (a^{\frac{2}{3}} +
    b^{\frac{2}{3}})^{\frac{3}{2}}$

    Least value $(a^{\frac{2}{3}} + b^{\frac{2}{3}})^{\frac{3}{2}},$ attained when $\tan^3 A = \frac{b}{a}$.

\item L.H.S. $= \sin^2 A + \cos^2 A + 4 + {\rm cosec}^2 A + \sec^2 A = 5 + \frac{1}{\sin^2 A\cos^2 A}$

    $\sin A\cos A \leq \frac{\sin^2 A + \cos^2 A}{2} = \frac{1}{2}$

    Least value $5 + 4 = 9,$ at $A = 45^\circ$.

\item Let $p = \sin A + \cos A \in (1, \sqrt{2}],$ $\sin A\cos A = \frac{p^2 - 1}{2}$

    $\tan A + \cot A = \frac{2}{p^2 - 1}, \sec A + {\rm cosec} A = \frac{2p}{p^2 - 1}$

    Sum $= p + \frac{2(p + 1)}{p^2 - 1} = p + \frac{2}{p - 1},$ which decreases on $(1, \sqrt{2}]$ as $(p - 1)(q - 1) < 1$ there

    Least value at $p = \sqrt{2}$: $\sqrt{2} + \frac{2}{\sqrt{2} - 1} = 2 + 3\sqrt{2}$ (at $A = 45^\circ$).

\item Let $p = \sin A + \cos A \leq \sqrt{2},$ $\sin A\cos A = \frac{p^2 - 1}{2}$

    (i) $(1 + \sin A)(1 + \cos A) = \frac{(p + 1)^2}{2} \leq \frac{(\sqrt{2} + 1)^2}{2} = \frac{3 + 2\sqrt{2}}{2}$

    (ii) $\frac{\sin A\cos A}{1 + \sin A + \cos A} = \frac{p^2 - 1}{2(p + 1)} = \frac{p - 1}{2} \leq \frac{\sqrt{2} - 1}{2}$

    Both greatest values occur at $A = 45^\circ$.

\item Let $t = \sin^2 A\cos^2 A,$ $0 < t \leq \frac{1}{4}$ for acute $A$.

    $\sin^6 A + \cos^6 A = 1 - 3t,$ so $\frac{1}{4} \leq 1 - 3t < 1$

    $\sin^8 A + \cos^8 A = (1 - 2t)^2 - 2t^2 = 1 - 4t + 2t^2 = 2(1 - t)^2 - 1$

    $1 - t \geq \frac{3}{4} \Rightarrow \sin^8 A + \cos^8 A \geq 2\cdot\frac{9}{16} - 1 = \frac{1}{8}$
\stopitemize
